\chapter{Droites, plans et systèmes linéaires}
\section*{Exercice 7.1 -- Équation de plan}
$2(x-1)-y+3(z-2)=0$, soit $2x-y+3z-8=0$. Substituer le point donne zéro ; les coefficients de $x,y,z$ restituent la normale.

\section*{Exercice 7.2 -- Intersection}
La substitution donne $t+1+t+2-t=3+t=6$, donc $t=3$. L'unique intersection est $(3,4,-1)$ ; la somme des coordonnées vaut bien $6$.

\section*{Exercice 7.3 -- Projection sur un plan}
Le défaut vaut $6$, la norme au carré de $(1,1,1)$ vaut $3$. Retirer $2(1,1,1)$ donne $H=(-1,0,1)$. La distance vaut $6/\sqrt3=2\sqrt3$ ; la somme des coordonnées de $H$ est nulle.

\section*{Exercice 7.4 -- Intersection de plans}
La différence impose $2y=-2$, donc $y=-1$ ; la première équation donne $x+z=2$. En posant $z=t$, on obtient la droite $(2-t,-1,t)$, de directeur $(-1,0,1)$.

\section*{Exercice 7.5 -- Paramètres indépendants}
L'égalité des coordonnées impose $s=0$, puis $t=1$. Elles se coupent en $(1,0,0)$. Exiger $s=t$ aurait artificiellement rendu le système incompatible et conduit à une conclusion fausse.

\section*{Exercice 7.6 -- Dégénérescence}
La première égalité est satisfaite par tout point de l'espace ; la seconde par aucun. Aucune ne définit un plan. La condition « normale non nulle » ne doit pas être omise.
